\section{Guarded operations on represented quantities}
\label{sec:guarded}

Approximate names support algebraic operations only when error propagation is
quantified.  Operations with singularities require explicit guards.  This section
states the generic reciprocal theorem and the broader interface discipline it
illustrates.

\subsection{Regular rational names}

A convenient real presentation is a rational evaluator \(a(k)\) with a fixed
coherence estimate, for example
\[
  \abs{a(k)-a(\ell)}\le2^{-k}+2^{-\ell}.
\]
Equivalent regular-name conventions are possible.  The important point is that a
finite output demand gives a primitive-recursive input precision through an error
budget.  Addition uses the sum of errors.  Multiplication additionally needs a bound
on the magnitudes of the factors.  Polynomial evaluation uses a finite telescoping
estimate or Horner recurrence with an explicit coefficient/value bound.

\subsection{Nonzero is not apartness}

Constructively, the proposition \(x\ne0\) is negative information.  It does not supply
a rational \(\lambda>0\) with \(\abs{x}\ge\lambda\).  A uniform operation that turned
an arbitrary represented nonzero real into such a \(\lambda\) would decide information
not continuously available from a Cauchy name.  Therefore reciprocal is a guarded
operation.

\begin{definition}[Apartness witness]
An apartness witness for a represented real \(x\) is positive rational data
\(\lambda>0\) together with a certificate \(\abs{x}\ge\lambda\).
\end{definition}

\begin{theorem}[Guarded reciprocal]
\label{thm:guarded-reciprocal}
Assume \(\abs{x}\ge\lambda>0\) and a rational approximation \(q\) satisfies
\(\abs{x-q}\le\eps\le\lambda/2\).  Then
\[
  \abs{q}\ge\lambda/2
\]
and
\[
  \left|\frac1x-\frac1q\right|
  \le\frac{2\eps}{\lambda^2}.
\]
Consequently a regular name for \(x\), together with \(\lambda\), yields a
primitive-recursive regular name for \(1/x\).
\end{theorem}

\begin{proof}
The reverse triangle inequality gives
\[
  \abs{q}\ge\abs{x}-\abs{x-q}
            \ge\lambda-\eps\ge\lambda/2.
\]
Therefore
\[
 \left|\frac1x-\frac1q\right|
 =\frac{\abs{x-q}}{\abs{x}\abs{q}}
 \le\frac{\eps}{\lambda(\lambda/2)}
 =\frac{2\eps}{\lambda^2}.
\]
All quantities may be replaced by integer numerator inequalities after clearing
positive denominators.  The object proof uses only ordered-ring arithmetic and the
apartness certificate.
\end{proof}

To obtain output error at most \(2^{-p}\), it suffices to request an input approximation
with
\[
  \eps\le\min\left(\frac\lambda2,
                        \frac{\lambda^2}{2^{p+1}}\right).
\]
A dyadic lower approximation to the right-hand side gives a primitive-recursive
precision transformer.

\subsection{Totalized guarded functions}

Every primitive-recursive symbol must be total.  Define a totalized reciprocal
approximator that returns a default value when the numerical guard check fails.  Its
contract has the form
\[
  \operatorname{Apart}(x,\lambda)
  \longrightarrow
  \operatorname{RecipCorrect}(x,\lambda,\operatorname{recip}(x,\lambda,p),p).
\]
The behavior outside the guard is irrelevant.  This is preferable to pretending the
operation is partial or hiding the guard in a semantic type class.

The same pattern applies to:

\begin{itemize}
\item division by a determinant or pivot;
\item normalization of a nonzero vector;
\item logarithms on a positive interval;
\item contour integrals away from poles;
\item projectors separated by a spectral gap; and
\item local inverse functions with a derivative lower bound.
\end{itemize}

\subsection{Pole clearing and algebraic reformulation}

Sometimes the guard can be avoided by changing the observable.  Instead of proving an
identity containing \(1/x\), multiply through by a power of \(x\) and prove a
polynomial identity.  Instead of selecting a normalized eigenvector, prove a statement
about the rank-one projector after clearing a norm denominator, or about the invariant
subspace.  Monic polynomial presentations similarly avoid division by a leading
coefficient.

This is not merely a workaround.  In a weak constructive theory, polynomial and
finite-algebra certificates are often the intrinsic content of a theorem, while a
normalized representative is an additional choice requiring apartness.

\subsection{Branches and exact discrete data}

Square roots, arguments, logarithms, and Puiseux leading terms combine continuous
approximation with discrete branch data.  A branch interface must include one of:

\begin{enumerate}[label=(\roman*),leftmargin=2.4em]
\item an exact algebraic presentation that determines the branch;
\item a rational region separated from every branch boundary;
\item overlapping relational branches with a stable downstream observable; or
\item an explicit boundary assumption supplying rank, sign, or monodromy data.
\end{enumerate}

An arbitrary Cauchy name does not generally determine exact discrete information by a
uniform finite procedure.  The obstruction is addressed explicitly in
\cref{sec:barriers}; the correct response is to strengthen the presentation, not to
hide an undecidable comparison inside a selector.
